% FREUID Challenge 2026 — Technical Report
% Compile: pdflatex freuid_technical_report.tex && bibtex freuid_technical_report && pdflatex x2
% (or: latexmk -pdf freuid_technical_report.tex)

\documentclass[11pt,a4paper]{article}

\usepackage[T1]{fontenc}
\usepackage[utf8]{inputenc}
\usepackage{lmodern}
\usepackage{microtype}
\usepackage{geometry}
\usepackage{graphicx}
\usepackage{booktabs}
\usepackage{amsmath}
\usepackage{hyperref}
\usepackage{xcolor}
\usepackage{enumitem}

\geometry{margin=2.5cm}
\hypersetup{
  colorlinks=true,
  linkcolor=blue!60!black,
  citecolor=blue!60!black,
  urlcolor=blue!60!black,
}

\newcommand{\teamname}{[TODO: Team Name]}
\newcommand{\kaggleteam}{[TODO: Kaggle Team Name]}
\newcommand{\contactemail}{[TODO: corresponding@email.com]}
\newcommand{\reportdate}{\today}

\title{%
  \textbf{The FREUID Challenge 2026}\\[0.4em]
  \large Technical Report --- \teamname
}
\author{%
  [TODO: Author Name(s)]\\
  \small
  [TODO: Affiliation]\\
  \texttt{\contactemail}
}
\date{Report date: \reportdate}

\begin{document}
\maketitle

\begin{abstract}
We detect fraudulent identity documents (physical manipulations, GenAI-generated
digital edits, and print-and-capture attacks) with a single fine-tuned ConvNeXt-V2-large
classifier. The key finding driving our result was that FREUID Score is dominated by
input resolution and by preserving native document aspect ratio: replacing square-padded
low-resolution inputs with an aspect-preserving rectangular resize and progressively
increasing training/inference resolution improved our public FREUID Score by roughly
two orders of magnitude (square 384px: 0.198 $\rightarrow$ rect 832px, no TTA: 0.00465),
without any architectural ensembling. We additionally applied capture-style
augmentation (perspective jitter, re-compression, sensor noise, tone/contrast jitter)
to close the domain gap between the training set (69{,}332 digital vs.\ only 20
physically captured images) and a test distribution that emphasizes print-and-capture
documents. Diverse-architecture ensembling (a normalizer-free NFNet) and multi-seed
averaging of the same recipe were both tried and both \emph{hurt} the public
leaderboard score once validated, despite promising offline agreement statistics; we
report these as negative results. Our selected submission scores \textbf{0.00465}
(FREUID Score, lower is better) on the public leaderboard. Code, the frozen model
checkpoint, and a network-isolated Docker inference sandbox are provided for full
reproducibility.
\end{abstract}

\noindent\textbf{Competition:} The FREUID Challenge 2026 (IJCAI-ECAI 2026)~\cite{freuid2026}\\
\textbf{Kaggle team:} \kaggleteam\\
\textbf{Code:} \url{https://github.com/[TODO: your-org]/[TODO: your-repo]} (commit \texttt{[TODO: SHA]})\\

\section{Introduction}
\label{sec:intro}

Identity document fraud detection must generalize across three qualitatively
different attack surfaces: (i) physical manipulation of a genuine document
(reprints, overlays, tampered security features), (ii) fully digital,
GenAI-assisted edits (text/face/layout changes), and (iii) print-and-capture
attacks, where a manipulated or synthetic document is physically printed and
re-photographed, destroying many of the pixel-level artifacts a purely digital
detector would rely on. The official metric, the FREUID Score, combines AuDET
with APCER at a fixed 1\% BPCER operating point (lower is better).

Our approach is intentionally simple: a single global-context vision
backbone, fine-tuned end-to-end at high, aspect-preserving input resolution,
with augmentation designed to simulate the print-and-capture domain gap
present in the released training data. We favored this over a
multi-branch/fusion architecture (global + patch-anomaly + OCR/layout +
physical-artifact branches) after observing that resolution and aspect-ratio
handling alone closed most of the gap to a strong baseline, and that adding a
second, architecturally diverse but weaker branch strictly hurt validated
leaderboard performance (Section~\ref{sec:results}).

\section{Method}
\label{sec:method}

\subsection{Overview}
The pipeline is: (1) decode the document image with OpenCV
(libjpeg-turbo), falling back to Pillow only if OpenCV cannot read the file;
(2) resize with an aspect-preserving rectangle (no padding); (3) run the
image through a fine-tuned ConvNeXt-V2-large backbone with a linear binary
head; (4) apply a sigmoid to obtain a fraud score in $[0,1]$. No
post-hoc score calibration, temperature scaling, or thresholding is applied —
the raw sigmoid output is submitted directly, matching the leaderboard's
numeric semantics (higher = more confident the document is fraudulent).

\subsection{Model architecture}
Backbone: \texttt{convnextv2\_large}~\cite{convnextv2} (timm~\cite{timm}), 1536-dim pooled
feature output, average pooling, \textbf{no frozen layers} (full end-to-end
fine-tune). Head: \texttt{Dropout(p=0.1)} $\rightarrow$ \texttt{Linear(1536,
1)} producing a single logit. Total forward pass is a single backbone
invocation per image; no test-time augmentation, no ensembling in the
selected submission. Inference batch size is 1 (see
Section~\ref{sec:inference} for why this does not affect numerics here).

\subsection{Training objective and procedure}
We fine-tune from ConvNeXt-V2's public ImageNet-pretrained weights (loaded
once at training time only — the released checkpoint we submit contains the
fully fine-tuned weights and requires no further downloads). Loss: focal loss
($\gamma=2.0$, $\alpha=0.75$) to handle the label imbalance
(40{,}005 genuine vs.\ 29{,}347 fraudulent in the training set). Optimizer:
AdamW with two learning-rate groups (backbone $1\times10^{-5}$, head
$1\times10^{-4}$), weight decay $1\times10^{-4}$, gradient clipping, and a
\texttt{OneCycleLR} schedule with linear warmup. We train for a single epoch
over the \emph{entire} labeled training set (no held-out fold, i.e.\
\texttt{--fit\_all}) once the recipe was locked in via earlier
fold-based experimentation, in order to maximize the data available to the
final model. Effective batch size 8 (\texttt{batch\_size=1},
\texttt{accumulate\_grad=8}) — required at our largest input resolution to
fit a 16\,GB GPU.

\textbf{Augmentation.} Standard augmentations: horizontal flip, affine
(translate/scale/rotate), motion/Gaussian blur, Gaussian noise,
brightness/contrast/hue/saturation jitter, JPEG re-compression, coarse
dropout. \emph{Capture-style} augmentation (enabled via
\texttt{--strong\_capture\_aug}) additionally applies: perspective jitter
(simulating an off-axis phone/scanner capture), a random choice of
downscale-then-upscale or aggressive JPEG re-encoding (simulating
re-photographing a document), sensor-realistic noise (Gaussian or ISO noise),
and gamma/CLAHE/brightness-contrast jitter (simulating varying capture
lighting). This was motivated directly by the training-set composition
(Section~\ref{sec:data}).

\subsection{Score calibration}
None applied. The sigmoid output of the fine-tuned head is submitted as-is at
both public and (planned) private evaluation; this matches how the public
leaderboard score was obtained, so no train/test mismatch is introduced by a
separate calibration step.

\section{Data}
\label{sec:data}

\subsection{FREUID training set}
69{,}352 labeled images across five document types/classes (\texttt{EGYPT/DL},
\texttt{GUINEA/DL}, \texttt{BENIN/DL}, \texttt{MOZAMBIQUE/DL},
\texttt{MAURITIUS/ID}), label distribution 40{,}005 genuine / 29{,}347
fraudulent. Critically, the \texttt{is\_digital} metadata field shows only
\textbf{20} images are physically captured (recaptured) versus 69{,}332
purely digital images — while the held-out/public test distribution
emphasizes print-and-capture documents much more heavily. This digital
$\rightarrow$ captured domain gap, not document type, is the dominant
generalization challenge in this competition (see the discussion log,
\texttt{discussion.txt}, in the repository for the full diagnostic trail that
led to this conclusion). We used capture-style augmentation
(Section~\ref{sec:method}) as our primary mitigation; a follow-up experiment
oversampling the 20 real captured rows within each training epoch was run
late in development — \textbf{[TODO: report its validated public-LB outcome
here once confirmed; see repository memory/discussion log for the
in-progress result at report-drafting time]}.

\subsection{External data and pre-trained models}
No external data beyond FREUID. The only external resource used is the
publicly available, permissively licensed ImageNet-pretrained
\texttt{convnextv2\_large} checkpoint distributed via \texttt{timm}
(downloaded once during training only). The released, fully fine-tuned
inference checkpoint requires \emph{no} network access and is included in
this repository (see \texttt{models/README.md}).

\subsection{Validation strategy}
Early architecture/resolution/augmentation decisions were validated with a
group-stratified out-of-distribution-by-document-type fold
(\texttt{StratifiedGroupKFold} on \texttt{type}). This fold strategy was
useful for catching gross regressions but \textbf{underestimated the
public-leaderboard digital/captured domain gap} (Section~\ref{sec:data}):
capture-style augmentation that looked neutral-to-harmful on the OOD-by-type
fold in fact improved the public leaderboard substantially, so later
development decisions (input resolution, inference scale, capture
augmentation strength, ensembling) were validated directly against the
public leaderboard rather than the offline fold. Private test labels were,
by construction, never available during model selection.

\section{Inference}
\label{sec:inference}
Test-time preprocessing matches training preprocessing except augmentation is
disabled: aspect-preserving rectangular resize to height 736 (width $=
\mathrm{round}(736 \times 1.585 / 32) \times 32 = 1152$), OpenCV decode with a
Pillow fallback, ImageNet normalization. \textbf{No test-time augmentation} is
applied in the selected submission (an earlier flip-averaged TTA variant was
strictly worse once the decode path was corrected — see
Section~\ref{sec:results}). Inference resolution (736) is deliberately
\emph{below} training resolution (832, $\approx 0.88\times$); this scale was
found empirically to sit at the minimum of a noisy but clearly unimodal
inference-scale response curve (704 / 736 / 768 / 800 / 832 all evaluated).
Batch size is 1; since ConvNeXt-V2 uses LayerNorm rather than BatchNorm,
batch size does not affect numerics at eval time, so this choice is purely
about matching the exact validated configuration and about being able to
attribute a decode failure to a single image in the Docker sandbox (see
\texttt{docker/README.md}). The exact command used to produce the selected
submission via the CSV workflow, and the equivalent Docker entrypoint used
for the private/final evaluation, are both documented in the repository's
top-level \texttt{README.md} and \texttt{docker/README.md}.

Hardware: single 16\,GB GPU (16{,}376\,MiB), WSL2/Linux, 54\,GB RAM + 64\,GB
swap. Scoring the 7{,}821 publicly-scored rows takes roughly 15--18 minutes
at scale 736 without TTA. \textbf{[TODO: measure and report CPU-only
wall-clock, and private-set wall-clock once the private images and their
count are available on/after July 13.]}

\section{Results}
\label{sec:results}

\begin{table}[h]
  \centering
  \caption{Selected leaderboard results (lower FREUID Score is better).}
  \label{tab:results}
  \begin{tabular}{lcc}
    \toprule
    Configuration & FREUID $\downarrow$ & Notes \\
    \midrule
    Square-pad, 384px                          & 0.198    & Early baseline \\
    Rect (aspect-preserving), 448px             & 0.16050  & \\
    Rect, 512px                                 & 0.15403  & \\
    Rect, 640px                                 & 0.05149  & Step-change vs.\ 512px \\
    Rect, 768px                                 & 0.02342  & \\
    Rect, 832px, infer @768, TTA                & 0.00515  & \\
    Rect, 832px, infer @736, TTA                & 0.00495  & Previous champion \\
    Rect, 832px, infer @736, \textbf{no TTA}    & \textbf{0.00465} & \textbf{Selected submission} \\
    \midrule
    NFNet (diverse arch.) standalone @736       & 0.11438  & Weaker; any ensemble weight on it hurt \\
    ConvNeXt/NFNet ensemble (best weight)        & 0.00538  & Still worse than ConvNeXt alone \\
    Two-seed ConvNeXt average (same recipe)      & 0.00685  & Hurt despite high pairwise agreement ($r{=}0.936$) \\
    \bottomrule
  \end{tabular}
\end{table}

Private leaderboard: \textbf{[TODO: fill in after the private test images are
released (July 13, 2026) and the final submission is scored, per the code
freeze in Section~\ref{sec:repro}]}.

\subsection{Negative results}
We highlight two results that looked promising by offline proxy metrics but
did not improve the validated (public leaderboard) score, since we believe
they are informative for other teams:
\begin{itemize}[noitemsep]
  \item \textbf{Diverse-architecture ensembling.} A normalizer-free NFNet
  (\texttt{dm\_nfnet\_f1}), trained with an otherwise identical recipe,
  reached a competitive training loss but scored 0.11438 standalone on the
  public leaderboard — over $20\times$ worse than the ConvNeXt-V2 champion.
  Every non-zero ensemble weight on the NFNet branch monotonically
  \emph{increased} (worsened) the FREUID Score. We conclude the public test
  distribution (skewed toward print-and-capture images) is far outside this
  architecture's effective generalization, despite good in-distribution
  training loss.
  \item \textbf{Multi-seed averaging.} Two independently-seeded runs of the
  identical champion recipe agreed strongly offline (Pearson $r = 0.936$,
  95.0\% sign agreement at the 0.5 threshold on the 7{,}821 publicly-scored
  rows) — a much stronger agreement signature than the NFNet pair, and the
  textbook precondition for variance-reducing ensembling. Nonetheless, both
  a 50/50 average (0.00685) and small-weight blends (5\%: 0.00507, 10\%:
  0.00510) scored \emph{worse} than the single champion checkpoint alone
  (0.00465), and the second seed's standalone score (0.03906) turned out to
  be far weaker than the first despite comparable training loss. We did not
  find a reliable offline signal that distinguished this failure case from a
  genuinely beneficial ensemble ahead of time.
\end{itemize}

\section{Reproducibility}
\label{sec:repro}
\begin{itemize}[noitemsep]
  \item \textbf{Repository:} \url{https://github.com/[TODO: your-org]/[TODO: your-repo]},
  commit \texttt{[TODO: 40-character SHA, frozen at or after the July 13
  private-test release per the competition's code-freeze policy]}.
  \item \textbf{Docker:} `docker build -f docker/Dockerfile -t freuid-repro:local .`
  (build context = repository root); weights are baked into the image from
  \texttt{models/convnext\_full\_drop01\_rect832\_capaug\_e1.pth} (SHA-256 in
  \texttt{models/README.md}) — no runtime downloads. Run with
  \texttt{docker run --network none -v /data:/data:ro -v
  /out:/submissions freuid-repro:local}. See \texttt{docker/README.md} for
  the full contract and environment-variable configuration.
  \item \textbf{Dependencies:} Python 3.11; see \texttt{docker/requirements.txt}
  for pinned, broadly-installable version ranges, and the top-level
  \texttt{README.md} for the exact development-environment versions
  (\texttt{torch 2.11.0+cu126}, \texttt{timm 1.0.27}, \texttt{opencv-python
  4.13.0}, \texttt{albumentations 2.0.8}).
  \item \textbf{Runtime:} single 16\,GB GPU recommended; CPU fallback is
  supported but not benchmarked in this draft --- \textbf{[TODO: fill in
  CPU wall-clock before the final submission]}.
  \item \textbf{Randomness:} training used a fixed seed (42) via
  \texttt{src/utils/seed.py::seed\_everything}; inference is deterministic
  given the frozen checkpoint (no dropout/augmentation at eval time, no TTA
  in the selected submission).
  \item \textbf{Network:} inference requires \emph{no} network access
  (\texttt{docker run --network none}); network is only used at
  \texttt{docker build} time to install Python dependencies.
\end{itemize}

\section{Team and contributions}
\textbf{[TODO: list team members and roles — modeling, data, engineering,
writing.]}

\section*{Acknowledgments}
\textbf{[TODO: optional funding/compute acknowledgments, organizer thanks.]}

\bibliographystyle{plain}
\bibliography{references}

\appendix
\section{Hyperparameters}
\label{app:hparams}
\begin{table}[h]
  \centering
  \small
  \begin{tabular}{ll}
    \toprule
    Hyperparameter & Value \\
    \midrule
    Backbone & \texttt{convnextv2\_large} (timm), full fine-tune \\
    Training image size & 832 (rect, aspect 1.585 $\Rightarrow$ 832$\times$1312) \\
    Inference image size & 736 (rect, aspect 1.585 $\Rightarrow$ 736$\times$1152) \\
    Batch size (train) & 1 (accumulate\_grad = 8, effective 8) \\
    Batch size (inference) & 1 \\
    Learning rate (head) & $1\times10^{-4}$ \\
    Learning rate (backbone) & $1\times10^{-5}$ \\
    Weight decay & $1\times10^{-4}$ \\
    Head dropout & 0.1 \\
    LR schedule & OneCycleLR, linear warmup \\
    Loss & Focal loss ($\gamma=2.0$, $\alpha=0.75$) \\
    Epochs & 1 (full training set, no held-out fold) \\
    Seed & 42 \\
    Test-time augmentation & None (selected submission) \\
    \bottomrule
  \end{tabular}
\end{table}

\section{Additional ablations}
See Table~\ref{tab:results} and the negative-results discussion in
Section~\ref{sec:results}. The full, chronological experimental log
(including every intermediate resolution/scale/augmentation configuration
tried) is preserved in \texttt{discussion.txt} in the repository for
transparency.

\end{document}
